\section*{Chapter \ref{ch:dv:american-options-and-early-exercise} --- American Options and Early Exercise}

\begin{solution}{exo:dv:american-options-and-early-exercise:1}
$6.0902-5.5735=0.5167$, 9.3\% of the European value.
\end{solution}

\begin{solution}{exo:dv:american-options-and-early-exercise:2}
$\gamma=2\times0.05/0.04=2.5$, $S^*=100\times2.5/3.5=71.43$; at 100 the put is worth
$28.57\times(100/71.43)^{-2.5}=12.32$.
\end{solution}

\begin{solution}{exo:dv:american-options-and-early-exercise:3}
Below the interest on the strike over the 60 days, $50(1-e^{-0.05\times60/365})=0.41$:
exercise then loses at least the interest.
\end{solution}

\begin{solution}{exo:dv:american-options-and-early-exercise:4}
Exercising kills the option's remaining value, which grows with volatility, and earns
the interest on the strike, which grows with the rate. A borrow fee acts like a
dividend yield: it lowers the forward, raises the value of waiting with a put and so makes
its early exercise less likely; on a call it can make early exercise worthwhile.
\end{solution}

\begin{solution}{exo:dv:american-options-and-early-exercise:5}
With one date at expiry the holder has no choice, which is the European option. Adding
a date adds a choice: the holder can always ignore it, so the value cannot fall
(the new supremum is over a larger set of stopping times).
\end{solution}

\begin{solution}{exo:dv:american-options-and-early-exercise:6}
The deep in-the-money put at 8\%, error $-0.064$: there the premium is large and the
option is exercised soon, while the approximation's single power law is fitted to a
perpetual shape and misprices the near boundary.
\end{solution}

\begin{solution}{exo:dv:american-options-and-early-exercise:7}
0.932, against an interest on the strike of 0.785: with 90 days left the call's put
component is worth 0.15, and the dividend must pay for it too.
\end{solution}

\begin{solution}{exo:dv:american-options-and-early-exercise:8}
The European formula has no \omterm{def:dv:american-options-and-early-exercise:premium}{early-exercise premium}, so it needs a higher volatility
(21.4\%) to reach the American price. Priced with an American model the put is at 20\%,
like the European options. The ``richness'' is the premium, not volatility.
\end{solution}

\begin{solution}{pb:dv:american-options-and-early-exercise:1}
\textbf{1.} 20.30.
\textbf{2.} 19.56.
\textbf{3.} Yes: 0.74 a share better.
\textbf{4.} $80(1-e^{-0.04\times29/365})=0.254$.
\textbf{5.} 0.255: the call is so deep in the money that its time value is the interest on
the strike plus a put worth a tenth of a cent.
\textbf{6.} $0.7446\times100\times10\,000$: USD~744\,600.
\textbf{7.} 45\% of it: USD~335\,000.
\textbf{8.} The writers who are not assigned: the clearing house assigns exercise notices
randomly to clearing members, which assign them to their customers; a writer left
unassigned keeps a call worth 19.56 that it sold on a share worth 20.30.
\textbf{9.} It does not know how many holders will exercise, nor which writers will be
assigned: the transfer is an expectation, not a hedgeable cash flow.
\textbf{10.} It must deliver shares at 80; if it held shares as a hedge, it delivers them
and receives 80 but does not receive the dividend on them.
\textbf{11.} At or slightly below 20.30: nobody should pay more than the exercise value, and
holders who can exercise should not sell below it.
\textbf{12.} About 19.56.
\textbf{13.} Invested $100.30-20.30=80$; after the ex-date it holds the dividend 1.00, a
share at 99.30 and a short call worth 19.56: $1.00+99.30-19.56-80=+0.74$.
\textbf{14.} Zero, less costs: the share is delivered at 80, the amount invested.
\textbf{15.} By the expected failure rate and its uncertainty, knowing that the gain is paid
only on the calls not assigned, and that the whole market plays the same trade.
\textbf{16.} Inattention, holding through a broker without instructions, and the belief that
options are exercised automatically before expiry.
\textbf{17.} Exercising on the client's behalf needs the client's consent and funding (the
client must pay the strike); brokers alert rather than exercise.
\textbf{18.} As if every holder exercises optimally (the conservative value), releasing the gain
on the calls that are not assigned.
\textbf{19.} A threshold dividend of 0.255; USD~335\,000 transferred to writers.
\textbf{20.} When a dividend before expiry is worth more than the interest on the strike over
the remaining life plus the call's put component (or when a borrow fee makes owning the
share valuable).
\end{solution}

\begin{solution}{iq:dv:american-options-and-early-exercise:1}
A call: only just before an ex-date, when the dividend exceeds the interest on the
strike plus the lost optionality (or on a hard-to-borrow share). A put: yes, when deep in the
money with positive rates, since exercising earns the interest on the strike.
\iqlookfor{the interest-on-strike argument both ways.}
\end{solution}

\begin{solution}{iq:dv:american-options-and-early-exercise:2}
The holder chooses when to stop, using only information available at the time, and will
choose the best rule: the value is the supremum over stopping times of the discounted
expected payoff under the pricing measure. A tree computes it backwards, taking at each node
the maximum of the exercise value and the discounted expected value of the next step (the
Snell envelope, One Quant Book 4, chapter 10).
\iqlookfor{optimal stopping, and backward induction as its algorithm.}
\end{solution}

\begin{solution}{iq:dv:american-options-and-early-exercise:3}
Time-independent pricing equation $\tfrac12\sigma^2S^2V^{\prime\prime}+rSV^{\prime}-rV=0$; solutions $S$ and
$S^{-\gamma}$ with $\gamma=2r/\sigma^2$; the put keeps $AS^{-\gamma}$; value matching and
smooth pasting at $S^*$ give $S^*=K\gamma/(1+\gamma)$.
\iqlookfor{the two boundary conditions.}
\end{solution}

\begin{solution}{iq:dv:american-options-and-early-exercise:4}
Holders should exercise tonight; some will not. Expect assignment on most but not all of the
position, so the stock hedge must be sized for delivery (hold the shares) while the dividend
on shares that are called away is lost; the calls not assigned become worth less than their
exercise value tomorrow, a gain that should not be counted on.
\iqlookfor{assignment uncertainty and its effect on the hedge.}
\end{solution}

\begin{solution}{iq:dv:american-options-and-early-exercise:5}
A closed-form approximation (quadratic or Bjerksund--Stensland) or a precomputed grid in
the quoting loop; a tree or a finite-difference grid, with dividends, for risk and end of
day. Check the approximation against the accurate pricer on the whole range of moneyness,
maturity and rates, especially deep in the money and near ex-dates, and fall back to the
grid where it fails.
\iqlookfor{speed where it is needed, accuracy where it matters, and a benchmark.}
\end{solution}

\begin{solution}{iq:dv:american-options-and-early-exercise:6}
Invert the American pricer (tree or grid) with the dividends and borrow, or subtract an
estimate of the \omterm{def:dv:american-options-and-early-exercise:premium}{early-exercise premium} and invert the European formula. With the European
formula the premium is read as volatility: deep in-the-money puts and calls before
dividends show volatilities that are too high, a fake smile.
\iqlookfor{the premium mistaken for volatility.}
\end{solution}
